\documentclass[11pt,a4paper]{article}
\usepackage[margin=27mm]{geometry}
\usepackage[T1]{fontenc}\usepackage{lmodern}
\usepackage{amsmath,amssymb,amsthm,mathtools,booktabs,microtype}
\usepackage[colorlinks=true,linkcolor=blue,citecolor=blue,urlcolor=blue]{hyperref}
\newtheorem{theorem}{Theorem}
\newtheorem{lemma}[theorem]{Lemma}
\newtheorem{corollary}[theorem]{Corollary}
\theoremstyle{remark}\newtheorem{remark}[theorem]{Remark}
\DeclareMathOperator{\rank}{rank}
\DeclareMathOperator{\Tr}{Tr}
\newcommand{\C}{\mathbb C}\newcommand{\Z}{\mathbb Z}
\newcommand{\fl}[1]{\left\lfloor #1\right\rfloor}
\newcommand{\cl}[1]{\left\lceil #1\right\rceil}
\title{Removing the binary and ternary factors\\from cyclic Weyl rank uncertainty}
\author{Lluis Eriksson\\\small Independent researcher, Sweden}
\date{24 October 2026\\Working draft 2; domain clarifications}
\begin{document}\maketitle
\begin{abstract}
For a nonzero complex matrix in the cyclic Weyl basis, let $u_d(s)$ be the
minimum rank under a support budget of $s$ coefficients. We give exact
formulas reducing $u_d$ to the profile in the part of the dimension coprime
to $6$. Removing a factor $2^k$ uses a minimum over binary subgroup costs.
Removing a factor $3^k$ uses products of the three costs $1,2,3$, with rank
multipliers $3,2,1$. The lower bounds follow from residue-sector
compressions and central Fourier blocks; the ternary two-residue case uses
a triangular minor on two sectors of a directed three-cycle. Upper bounds
are attained by scalar cyclic projectors and binomials, combined with the
Chinese remainder theorem. In particular, in every even dimension whose
only prime factors are $2$ and $3$, the exact minimum is
$d/\max\{q:q\mid d,\ q\le s\}$. We obtain sharp uniform coefficient-erasure
budgets for arbitrary complex low-rank matrices. Residual profiles with
prime factors at least $5$, efficient recovery, and numerical stability
are not determined.
\end{abstract}

\section{Profiles and the main reductions}
Put $\omega_d=e^{2\pi i/d}$ and
\[
 X_de_j=e_{j+1\bmod d},\quad Z_de_j=\omega_d^je_j,\quad
 W_d(a,b)=X_d^aZ_d^b,\qquad (a,b)\in(\Z/d)^2.
\]
These $d^2$ matrices are orthogonal for the Hilbert--Schmidt inner product.
For $C\ne0$, let $s_W(C)$ count the nonzero coefficients in its unique Weyl
expansion. For every integer $s\ge1$, define
\[
 u_d(s)=\min\{\rank C:C\in M_d(\C),\ C\ne0,\ s_W(C)\le s\}.
\]
The minimum exists: the admissible set is nonempty and its ranks form a
nonempty subset of $\{1,\ldots,d\}$. Thus no compactness argument is needed.
The profile is nonincreasing, $u_d(1)=d$, and $u_d(s)=1$ for $s\ge d$,
by the diagonal rank-one projector. We use $u_1(s)=1$.

\begin{theorem}[Binary transfer]\label{thm:two}
For odd $m\ge1$, $k\ge0$ and $s\ge1$,
\[
 u_{2^km}(s)=
 \min_{\substack{0\le a\le k\\2^a\le s}}
       2^{k-a}u_m\!\left(\fl{s/2^a}\right).
\]
\end{theorem}
For ternary transfer define, for $3\nmid m$,
\[
 F_{k,m}(s)=
 \min_{\substack{r_1,\ldots,r_k\in\{1,2,3\}\\
                  r_1\cdots r_k\le s}}
 \left(\prod_{i=1}^k(4-r_i)\right)
 u_m\!\left(\fl{s/(r_1\cdots r_k)}\right).
\]
The empty tuple gives $F_{0,m}(s)=u_m(s)$.
\begin{theorem}[Ternary transfer]\label{thm:three}
For $3\nmid m$, $m\ge1$, $k\ge0$ and $s\ge1$,
\[
 u_{3^km}(s)=F_{k,m}(s).
\]
\end{theorem}
The unknown residual profiles are retained as functions, not replaced by
an assumed uncertainty conjecture. All coefficients are arbitrary complex
numbers, and sharpness means a support budget, not exactly $s$ terms.

\begin{corollary}[An even residual formula]\label{cor:residual}
Let $h=2^a3^b$ with $a\ge1$, $b\ge0$, and let $\gcd(m,6)=1$. Then
\[
 u_{hm}(s)=
 \min_{\substack{q\mid h\\q\le s}} \frac hq\,
                  u_m\!\left(\fl{s/q}\right).
\]
In particular, if $d=2^a3^b$ is even and
$Q_d(s)=\max\{q:q\mid d,\ q\le s\}$, then
\[
                         u_d(s)=d/Q_d(s).
\]
\end{corollary}

The scalar uncertainty profile in prime-power dimension is classical
\cite{dvb}; Meshulam's subgroup--quotient argument explicitly uses a
smallest occupied coset and a product constraint on support sizes
\cite[Proposition 1.3 and Section 2]{meshulam}. We credit that architecture.
The proposed operator contribution is the exact transfer of the rank
profile through the noncommuting cyclic representation, using the block
identities below. Earlier work in the author's programme covers low
prime-power sparsity and a binary seven-term range \cite{prior,seven}.
The present argument replaces the binary overlap classification and
treats the ternary profile in every sparsity range.

Yin studies the closely related Hermitian traceless rank-two problem in
all cyclic dimensions, together with phase-retrieval and equality
statements \cite{yin}. Our profile allows arbitrary complex matrices and
all rank bounds, but does not give the stronger positivity, trace, or
phase-retrieval classifications of that setting. The reference's
formalization claims have not been replayed here. These comparisons are
limited readings, not an exhaustive certificate of priority.
The cyclic basis is not being identified with a tensor-qubit Pauli basis.

\section{Compression identities and small Fourier bounds}
Left multiplication by a Weyl unitary translates the support, up to
nonzero phases, and preserves rank. Discrete Fourier conjugation
interchanges the two label coordinates, up to a sign and phases.

\begin{lemma}[Residue sectors]\label{lem:residue}
Let $d=pn$ for $p=2$ or $3$ and
$E_j=\operatorname{span}\{e_{j+pu}:u\in\Z/n\}$, $0\le j<p$.
For $a=pA+\delta$, $0\le\delta<p$, put
$j'=(j+\delta)\bmod p$ and $\kappa=\fl{(j+\delta)/p}$. Then
\[
 W_d(a,b)|_{E_j\to E_{j'}}
       =\omega_d^{bj}X_n^{A+\kappa}Z_n^b .
\]
Within a fixed first-coordinate residue class, distinct labels merge
precisely in cosets of $\langle(0,n)\rangle$.
An occupied class gives a nonzero block somewhere. A coset with $r$
occupied labels activates at least $p-r+1$ blocks, for $1\le r\le p$.
\end{lemma}
\begin{proof}
The output index on $e_{j+pu}$ is $j'+p(u+A+\kappa)$ modulo $d$,
and the phase is $\omega_d^{bj+ pbu}$. This is the displayed formula.
The compressed label is $(A\bmod n,b\bmod n)$ after factoring the common
carry. In a merging coset, write $b=\beta+n\ell$; its coefficient, after
a nonzero common factor, is
$\sum_{\ell=0}^{p-1}c_\ell e^{2\pi i\ell j/p}$.
Different cosets give different compressed Weyl labels. Thus an occupied
class cannot vanish in all blocks, by invertibility of this transform.

For $p=2$, one input activates both outputs and two inputs activate at
least one. For $p=3$, one input activates three outputs; two inputs cannot
have just one output, because the inverse of a single Fourier spike
has all three entries nonzero; three inputs activate at least one.
These are precisely the claimed bounds. No assumption on coefficient
signs or algebraicity is used.
\end{proof}

\begin{lemma}[Central blocks]\label{lem:central}
Let $d=p^2n$, $p=2$ or $3$, and suppose the support lies in $p(\Z/d)^2$.
The commuting unitaries $X_d^{pn},Z_d^{pn}$ are central on this support
algebra. Their $p^2$ joint eigenspaces have dimension $n$, are invariant
under $C$, and carry cyclic Weyl representations. Distinct cosets of
$H=\langle(pn,0),(0,pn)\rangle$ give distinct Weyl labels in every block.
Within an $H$-coset the coefficients transform by ordinary Fourier
transform on $(\Z/p)^2$, apart from nonzero scalar factors.
\end{lemma}
\begin{proof}
Let $\zeta_p=e^{2\pi i/p}$. For $\alpha,\beta\in\{0,\ldots,p-1\}$, the
joint eigenspace has the orthogonal basis
\[
 v_t^{\alpha,\beta}=\omega_d^{-p\beta t}
 \sum_{\ell=0}^{p-1}\zeta_p^{-\beta\ell}
 e_{\alpha+pt+pn\ell},\qquad 0\le t<n .
\]
Its vectors have norm $\sqrt p$; normalizing gives orthonormal bases.
The two eigenvalues are $\zeta_p^\beta$ and $\zeta_p^\alpha$.
The bases exhaust $\C^d$. Direct action, including wrap-around, gives
\begin{equation}\label{eq:central}
 W_d(pA,pB)|_{E_{\alpha,\beta}}
       =\omega_d^{p\beta A+p\alpha B} X_n^AZ_n^B .
\end{equation}
The factor $\omega_d^{p^2}$ is primitive of order $n$, also for arbitrary
$n$. Representatives $0\le A,B<n$ label the $n^2$ distinct compressed Weyls.

The labels in one coset are
$(pA+pn\epsilon,pB+pn\eta)$, $0\le\epsilon,\eta<p$.
By \eqref{eq:central} its compressed coefficient is
\[
 \omega_d^{p\beta A+p\alpha B}
 \sum_{\epsilon,\eta=0}^{p-1}
        c_{\epsilon,\eta}\zeta_p^{\beta\epsilon+\alpha\eta}.
\]
This proves both the Fourier assertion and injectivity between cosets.
The formula remains valid when $n=1$.
\end{proof}

Put $D(r)=2^{\fl{\log_2r}}$ for $r\ge1$, and define
\[
 \phi_9(r)=
 \min_{\substack{u,v\in\{1,2,3\}\\uv\le r}}(4-u)(4-v),
 \qquad 1\le r\le9.
\]
\begin{lemma}[Central Fourier support]\label{lem:fourier}
A nonzero vector on $(\Z/2)^2$ with $r$ occupied entries has at least
$4/D(r)$ nonzero Fourier coefficients. On $(\Z/3)^2$, the corresponding
lower bound is $\phi_9(r)$.
\end{lemma}
\begin{proof}
For four entries, $r=1$ gives four outputs. For $r=2,3$, a single output
is impossible because its inverse has all four entries nonzero. For $r=4$,
invertibility suffices. The bound at $r=3$ is deliberately not sharpened.

For nine entries, group by one coordinate. Let $v$ be the number of
occupied fibers and $u$ the smallest fiber support; then $uv\le r$.
Transform that fiber first: at least $4-u$ coordinate frequencies remain.
Each such nonzero vector along the other coordinate has support at most
$v$, and hence at least $4-v$ Fourier outputs by Lemma~\ref{lem:residue}'s
three-point argument. Thus there are at least $(4-u)(4-v)\ge\phi_9(r)$
outputs. This is the elementary coset argument credited above.
\end{proof}

\begin{lemma}[Two shifts on three sectors]\label{lem:cycle}
Suppose a matrix on three equal-sized input/output sectors has exactly
two possible block shifts. If one shift has $L$ nonzero blocks, each of
rank at least $\rho$, then the whole matrix has rank at least
$\min(L,2)\rho$.
\end{lemma}
\begin{proof}
Permute output sectors to place the chosen shift on the diagonal. The
other shift is a nonzero translation on $\Z/3$, a directed three-cycle.
Any two sectors induce at most one off-diagonal direction, so their
block submatrix is triangular after an ordering. Select a nonsingular
minor of full rank in each of its two diagonal blocks. Together these
give a triangular minor with nonzero determinant, even when the full
diagonal blocks are singular. If $L=1$, use that block alone.
\end{proof}

\section{Upper bounds and attainable costs}
For coprime $r,m$, the CRT unitary $e_j\mapsto e_{j\bmod r}\otimes
e_{j\bmod m}$ sends
\[
 X_{rm}\mapsto X_r\otimes X_m,\qquad
 Z_{rm}\mapsto Z_r^{m^{-1}\bmod r}\otimes Z_m^{r^{-1}\bmod m}.
\]
Consequently it bijectively identifies Weyl labels with pairs of labels.
Tensor products multiply ranks and support sizes. A factor of dimension
one is treated directly, without modular-inverse notation.

For $q\mid N$, let
\[
 P_q=\sum_{\ell=0}^{q-1}Z_N^{\ell N/q}.
\]
It equals $q$ on diagonal positions divisible by $q$, and zero elsewhere.
Thus it has support $q$ and rank $N/q$. This proves the upper bound in
Theorem~\ref{thm:two} by taking $N=2^k$, $q=2^a$ and tensoring a
matrix attaining the residual budget.

For ternary transfer, two entries $r_i=r_j=2$ in the product defining
$F_{k,m}$ may be replaced by $3,1$: cost drops from $4$ to $3$, the rank
multiplier drops from $4$ to $3$, and the residual budget increases.
Hence the minimum is attained with at most one $2$. It is the minimum
of the following two families, retaining only positive residual budgets:
\begin{align*}
 &3^{k-a}u_m(\fl{s/3^a}),&&0\le a\le k,\quad3^a\le s,\\
 &2\,3^{k-a-1}u_m(\fl{s/(2\,3^a)}),
       &&0\le a<k,\quad2\,3^a\le s .
\end{align*}
The first is attained by $P_{3^a}$ on $N=3^k$. For the second use
\[
 P_{3^a}\bigl(Z_N^{N/(3\cdot3^a)}-\zeta_3 I\bigr).
\]
The two exponent sets are disjoint modulo $N$, giving exactly $2\,3^a$
terms. Among the $N/3^a$ nonzero positions of $P_{3^a}$, one third
vanish in the second factor. Its rank is $2N/(3\cdot3^a)$.
CRT and a residual minimizer now prove the upper bound of
Theorem~\ref{thm:three}.

\section{Lower bounds by induction}
We prove the two transfer theorems separately by induction on $k$,
holding $m$ fixed. It suffices to treat actual support $s_0$, because
both proposed minima are nonincreasing in the budget.

\subsection{Binary transfer}
The base $k=0$ is the definition of $u_m$. If some coordinate has both
parities, exchange coordinates if necessary and take the smaller
first-coordinate class. Lemma~\ref{lem:residue} gives a nonzero block
of dimension $d/2$ with support at most $\fl{s_0/2}$. Induction bounds its
rank by the minimum over indices $1\le a'\le k$ with $2^{a'}\le s_0$ of
$2^{k-a'}u_m(\fl{s_0/2^{a'}})$. This is at least the desired minimum,
which additionally permits $a'=0$.

Otherwise translate the support to even labels. If $k=1$, $d=2m$,
the two invariant sectors carry
\[
 W_d(2A,2B)\mapsto \omega_m^{Bj}X_m^AZ_m^{2B},
 \qquad j=0,1.
\]
Multiplication by $2$ is invertible modulo $m$. The blocks are equivalent:
conjugation by $X_m^h$ with $-2h\equiv1\pmod m$ changes the first into
the second. Thus $\rank C\ge2u_m(s_0)$, the $a'=0$ candidate.
For $m=1$ this case has only one label and is immediate.
These are two equivalent representations, not two independently
variable algebra blocks.

If $k\ge2$, put $n=d/4$ and apply Lemma~\ref{lem:central}.
Let $t$ be the number of occupied central cosets and $r$ their minimum
occupancy, so $rt\le s_0$ and $1\le r\le4$.
At least $4/2^\ell$ blocks are nonzero, where $\ell=\fl{\log_2r}$.
Each has support at most $t$. Its rank is at least $u_n(t)$, and ranks
add over invariant blocks. For every index $a$ in the inductive formula
for $u_n(t)$, put $a'=a+\ell$. Then
\[
 a'\le k,\qquad 2^{a'}\le rt\le s_0,\qquad
 \fl{t/2^a}\le\fl{s_0/2^{a'}} .
\]
Its scaled rank candidate is therefore at least
$2^{k-a'}u_m(\fl{s_0/2^{a'}})$, and hence at least the desired minimum.
All arguments of the profiles are positive. Their definition for every
positive budget handles $t\ge n$ without an out-of-domain induction.
This proves the lower bound of Theorem~\ref{thm:two}.

\subsection{Ternary transfer}
The base is again $k=0$. Suppose some coordinate varies modulo $3$,
and make it the first. If all three first-coordinate residues occur,
a smallest class has at most $\fl{s_0/3}$ terms and supplies a nonzero
block. Thus
\[
 \rank C\ge F_{k-1,m}(\fl{s_0/3})\ge F_{k,m}(s_0).
\]
The last inequality follows by prepending $r=3$ to each eligible tuple.

If exactly two residues occur, choose a smaller class of size
$M\le\fl{s_0/2}$. Let it occupy $t$ cosets of $\langle(0,d/3)\rangle$,
and let $r\le3$ be the smallest coset size. Then $rt\le M$.
At least $4-r$ blocks are nonzero, each of rank at least $u_{d/3}(t)$.
By Lemma~\ref{lem:cycle}, $r=1$ or $2$ gives
\[
 \rank C\ge2F_{k-1,m}(\fl{s_0/2})\ge F_{k,m}(s_0),
\]
by monotonicity and prepending $r=2$. If $r=3$, necessarily $s_0\ge6$,
and a single nonzero block yields
\[
 \rank C\ge F_{k-1,m}(\fl{s_0/6})
          \ge F_{k-1,m}(\fl{s_0/3})\ge F_{k,m}(s_0).
\]
Thus neither existence nor rank additivity for arbitrary blocks is assumed.

It remains that both coordinates have fixed residues. Translate to $3G$.
When $k=1$, $3\nmid m$, the three invariant sectors carry
\[
 W_d(3A,3B)\mapsto\omega_m^{Bj}X_m^AZ_m^{3B},\qquad j=0,1,2 .
\]
These have distinct labels, and are conjugate by $X_m^h$ with
$-3h\equiv j\pmod m$. Their equal ranks sum to at least $3u_m(s_0)$,
the all-$r_i=1$ candidate. The case $m=1$ is immediate.

For $k\ge2$, put $n=d/9$. Lemmas~\ref{lem:central} and
\ref{lem:fourier} give nine invariant blocks. If $t$ cosets are occupied
and $r$ is their minimum occupancy, at least $\phi_9(r)$ blocks are nonzero.
Thus
\[
                   \rank C\ge\phi_9(r)F_{k-2,m}(t),\qquad rt\le s_0.
\]
Choose $u,v$ attaining $\phi_9(r)$, and a tuple attaining $F_{k-2,m}(t)$
with cost $w\le t$. The concatenated tuple is eligible for $F_{k,m}(s_0)$,
because $uv\le r$ and $uvw\le s_0$. Also
$\fl{t/w}\le\fl{s_0/(uvw)}$. Monotonicity of $u_m$ makes that candidate
at most $\phi_9(r)F_{k-2,m}(t)$, proving the required bound.
This completes Theorem~\ref{thm:three} in its entire stated domain.

\section{Removing both factors and the divisor law}
\begin{proof}[Proof of Corollary~\ref{cor:residual}]
Apply Theorem~\ref{thm:two} with odd residual $3^bm$, then
Theorem~\ref{thm:three}. Floor composition is exact for positive integer
costs. Candidates without a ternary cost $2$ are
\[
 2^{a-i}3^{b-j}u_m(\fl{s/(2^i3^j)}),\qquad
 0\le i\le a,\ 0\le j\le b .
\]
These are precisely the divisor candidates in the corollary.
A candidate with a ternary cost $2$ has cost $2^{i+1}3^j$ and multiplier
$2^{a-i+1}3^{b-j-1}$, with $j<b$.

If $i<a$, replace it by the divisor candidate at cost $2^{i+1}3^j$:
its multiplier is $2^{a-i-1}3^{b-j}$, smaller by a factor $3/4$,
and its residual budget is unchanged. If $i=a$, use instead the divisor
cost $q'=2^{a-1}3^{j+1}$, which is $3/4$ of the former cost.
Its multiplier is the same, $2\,3^{b-j-1}$, and its residual budget
is larger. Thus every additional candidate is dominated. For $m=1$
the residual profile is constant one, and minimizing $d/q$ chooses
the largest eligible divisor.
\end{proof}

For a purely ternary dimension $d=3^k$, the same proof also gives
\[
 u_d(s)=(4-c)3^{k-t-1}
 \quad\text{when}\quad c3^t\le s<(c+1)3^t,\quad c\in\{1,2\},\ s<d,
\]
and $u_d(s)=1$ for $s\ge d$. For a purely binary dimension it gives
$d/2^{\min(k,\fl{\log_2s})}$. The scalar profiles are recovered as
operator profiles here; no scalar theorem is silently used to justify
the noncommuting lower bound.

\begin{center}
\begin{tabular}{@{}rrrrr@{}}
\toprule
$d$ & budget $s$ & product bound $\cl{d/s}$ & exact $u_d(s)$ & subgroup cost\\
\midrule
54 & 8 & 7 & 9 & 6\\
72 & 11 & 7 & 8 & 9\\
108 & 8 & 14 & 18 & 6\\
108 & 17 & 7 & 9 & 12\\
\bottomrule
\end{tabular}
\end{center}
The elementary product bound follows from
$\|C\|_F^2=d\sum|c_g|^2$ and
$\|C\|_{\rm op}^2\le s\sum|c_g|^2$.
The divisor law is specific to the stated dimensions. For example,
in dimension $10$, a degree-three polynomial in $Z_{10}^2$ with three
distinct fifth roots as zeros has at most four Weyl terms and rank four.
This is below $10/\max\{q\mid10:q\le4\}=5$. There is no claim of a
divisor-step law for all composite dimensions.

\section{Sharp erasure budgets}
Let $d$ be even with only prime factors $2,3$. For an integer $R\in\{1,\ldots,d\}$, set
\[
 q_d(R)=\min\{q:q\mid d,\ qR\ge d\}.
\]
\begin{corollary}\label{cor:erase}
The minimum support of a nonzero complex matrix of rank at most $R$ is
$q_d(R)$. On the set of all complex matrices of rank at most $R$,
the largest universally safe number of omitted Weyl coefficients is
$q_d(2R)-1$ if $2R<d$, and zero if $2R\ge d$.
\end{corollary}
\begin{proof}
The first statement is the inverse of the divisor profile; the witness
$P_{q_d(R)}$ attains it. If two matrices of rank at most $R$ have the
same retained coefficients, their difference has rank at most $2R$
and support contained in the omissions. A nonzero difference needs at
least $q_d(2R)$ coefficients.

For failure at that number, use $P_q$, $q=q_d(2R)$, whose rank is
$d/q\le2R$. Split its nonzero diagonal positions into two groups of
size at most $R$. Put $A=P_q$ on the first group and $B=-P_q$ on the
second, zero elsewhere. Then $A-B=P_q$, while both ranks are at most $R$.
If $2R\ge d$, the same split of $I_d$ shows failure on erasing its single
Weyl label. No omissions retains a basis and is injective.
\end{proof}
For $d=108$ and $R=4$, the sharp budget is $17$, with failure at $18$.
The product bound alone guarantees $13$. This is exact noiseless
identifiability of arbitrary complex matrices. It does not establish a
recovery algorithm, conditioning estimate, or a positivity-constrained
density-matrix threshold.

\section{Checks, dependencies and status}
The proof is analytic for all stated dimensions. Exact auxiliary checks
compare root-of-unity monomial exponents, CRT label maps and the integer
minimum identities. Finite checks do not establish the universal
quantifiers. In particular, test residual functions used for minimum
bookkeeping are explicitly artificial monotone functions, not asserted
physical Weyl rank profiles.

The binary and ternary transfer arguments have received dependent
mathematical audits. A full assessment of draft~1 requested two domain
clarifications: integer rank budgets in Corollary~\ref{cor:erase}, and
admissible positive budgets in the binary induction. Both are explicit in
this revision. That assessment and its findings are preserved separately;
the revised artifact has not yet received its own assessment. An assessment
of the earlier binary-only PDF does not apply to this text.

The residual profile for dimensions with prime factors at least $5$ is
not determined in general. The two-shift proof uses the special directed
three-cycle on the sector labels; it is not asserted for an arbitrary
prime number of sectors. All sharpness assertions are budget statements,
not classifications of exact-support minimizers or of all equality cases.

\paragraph{AI and responsibility disclosure.}
The named author directed OpenAI Codex assistance in derivation, writing
and exact checks. GPT-6 Astra xhigh audits used shared filesystem and
declared reused task context. No runtime isolation, independent human
review, Lean formalization, certified worldwide originality or predicted
assessment score is claimed. Reports on other exact artifacts are
preserved separately and their grades are not transferred.

\begin{thebibliography}{9}
\bibitem{meshulam}
R. Meshulam, An uncertainty inequality for finite abelian groups,
\emph{European J. Combin.} \textbf{27} (2006), 63--67.
\url{https://arxiv.org/abs/math/0312407v1}.
\bibitem{dvb}
S. Delvaux and M. Van Barel, Rank-deficient submatrices of Fourier matrices,
\emph{Linear Algebra Appl.} \textbf{429} (2008), 1587--1605.
\href{https://doi.org/10.1016/j.laa.2008.04.043}{doi:10.1016/j.laa.2008.04.043}.
Author report TW470 (2006):
\url{https://www.cs.kuleuven.be/publicaties/rapporten/tw/TW470.pdf}.
\bibitem{prior}
L. Eriksson, \emph{Sparse cyclic Weyl operators in prime-power dimension:
an operator extension of the scalar uncertainty law through sparsity
$2p-1$}, AIRR.SCIENCE, ARR-2026-1KTM07SEEV968BCB, v1 (2026).
\href{https://airr.science/papers/ARR-2026-1KTM07SEEV968BCB/versions/v1/}{Version record}.
\bibitem{seven}
L. Eriksson, \emph{Sparse Weyl operators: a uniform binary rank theorem
and sharp erasure thresholds}, AIRR.SCIENCE,
ARR-2026-161TZ0CP488XTTBN, v1 (2026).
\href{https://airr.science/papers/ARR-2026-161TZ0CP488XTTBN/versions/v1/}{Version record}.
\bibitem{yin}
P. Yin, \emph{A Sharp Rank-Two Uncertainty Principle for Finite Weyl
Transforms}, public manuscript, commit
\texttt{39af48a5cf2e014a14f20e3d0ce836a430fea394},
\href{https://github.com/holdonyb/weyl-stft-paper/blob/39af48a5cf2e014a14f20e3d0ce836a430fea394/research/manuscript.tex}{pinned source},
accessed 24 October 2026.
\end{thebibliography}
\end{document}
